\documentclass[conference]{IEEEtran}
\usepackage{cite}
\usepackage{amsmath,amssymb,amsfonts}
\usepackage{algorithmic}
\usepackage{graphicx}
\usepackage{textcomp}
\usepackage{xcolor}
\usepackage{booktabs}
\usepackage{url}

\def\BibTeX{{\rm B\kern-.05em{\sc i\kern-.025em b}\kern-.08em
    T\kern-.1667em\lower.7ex\hbox{E}\kern-.125emX}}
\begin{document}

\title{OPTI-FAB: Latency-First Stream-Aware Inference with Confidence-Gated Early Exit for Real-Wafer Wafer Defect Classification}

\author{\IEEEauthorblockN{Hariharan M.}
\IEEEauthorblockA{\textit{Independent Researcher} \\
Email: hariharan@example.com}
\and
\IEEEauthorblockN{Asmitha M.}
\IEEEauthorblockA{\textit{Independent Researcher} \\
Email: asmitha@example.com}
}

\maketitle

\begin{abstract}
Automated visual inspection of semiconductor wafers using deep neural networks is critical for modern fab yield management. However, traditional deployment pipelines process images in a serialized capture-save-load-infer loop, where over 70\% of the end-to-end latency is spent on file I/O, OS buffering, and memory transfers. This latency bottleneck makes inline wafer rejection at transport speeds (e.g., 100~mm/s) practically unfeasible. In this work, we present OPTI-FAB, a latency-first edge-AI framework that reframes wafer inspection as a streaming-systems problem. OPTI-FAB processes inspection data as a continuous pixel stream using a circular sliding-window buffer, overlapping acquisition with inference. We introduce a confidence-gated decision mechanism utilizing normalized Shannon entropy to trigger early exits, allowing defect classification before full-frame capture. For borderline cases, we implement a dual-path architecture that routes uncertain frames to a Monte Carlo (MC) Dropout pipeline for epistemic uncertainty estimation. Optimized with TensorRT FP16, our model achieves a mean inference latency of 0.79~ms on an NVIDIA RTX 4050 Edge GPU, enabling a 63.9\% reduction in end-to-end pipeline latency (3.00~ms vs. 8.30~ms for file-based pipelines). OPTI-FAB delivers 90\% classification accuracy across eight wafer defect classes, proving that inline quality control at full wafer transport speeds is highly feasible.
\end{abstract}

\begin{IEEEkeywords}
Wafer inspection, Edge AI, Streaming inference, TensorRT, Early exit, Uncertainty estimation
\end{IEEEkeywords}

\section{Introduction}
Semiconductor manufacturing requires rapid, high-precision quality control to minimize yield loss. Wafer surface defect inspection, historically performed manually or via classical rule-based computer vision, has increasingly transitioned to deep learning-based object detection and classification models. Despite the high accuracy of these modern convolutional networks, their practical integration into production lines is often limited by strict system latency constraints.

Wafer transportation systems typically move wafers under scanning line-scan cameras at speeds of 100~mm/s or higher. To achieve true inline rejection---where a defective wafer is caught and diverted before exiting the inspection zone---the classification pipeline must complete within a tight window of a few milliseconds. 

Traditional visual inspection setups utilize a serialized processing pipeline:
\begin{equation}
\text{Capture} \rightarrow \text{Buffer} \rightarrow \text{Save File} \rightarrow \text{OS I/O} \rightarrow \text{Load} \rightarrow \text{Infer} \rightarrow \text{Decide}
\end{equation}
System-level profiling reveals that more than 70\% of the end-to-end latency in this pipeline is spent outside the neural network. Overhead is dominated by local storage writes, operating system file caches, disk I/O, and GPU host-to-device memory copies. Optimizing model weights or pruning channels alone cannot address this I/O bottleneck.

In this paper, we propose \textbf{OPTI-FAB}, a latency-first inspection system designed to bypass file-system serialization entirely. OPTI-FAB treats visual inspection as a continuous pixel stream:
\begin{equation}
\text{Camera Stream} \rightarrow \text{Circular Buffer} \leftrightarrow \text{Inference} \rightarrow \text{Gate} \rightarrow \text{Decide}
\end{equation}
By maintaining a circular sliding-window buffer, OPTI-FAB overlaps pixel ingestion with edge GPU inference. Crucially, rather than waiting for a full frame to be scanned and saved, OPTI-FAB runs inference on partially filled frames. Using a confidence-gated early exit, the system makes deterministic classification decisions once cumulative confidence exceeds a calculated threshold (typically at 50\% to 68\% of frame scanning completion).

To ensure robustness, OPTI-FAB incorporates a dual-path inference design. High-confidence deterministic predictions are routed through an ultra-fast TensorRT FP16 path. Borderline or noisy inputs are analyzed using Monte Carlo (MC) Dropout on a Keras backend to estimate epistemic uncertainty, allowing the system to flag highly uncertain wafers for manual operator review instead of risking false accepts.

\section{Methodology}

\subsection{Problem Formulation and Latency Budget}
Given a wafer transport mechanism operating at a constant velocity $v = 100\text{ mm/s}$ and a physical tolerance window $d = 0.3\text{ mm}$ for inline mechanical rejection gates, the maximum allowable end-to-end latency budget $T_{\text{budget}}$ is defined as:
\begin{equation}
T_{\text{budget}} = \frac{d}{v} = 3.0\text{ ms}
\end{equation}
If the pipeline latency exceeds 3.0~ms, the wafer will have traveled past the physical rejection mechanism, making inline sorting impossible.

\subsection{Circular Buffer Sliding-Window}
A camera scanner feeds grayscale pixel rows line-by-line. OPTI-FAB maintains a fixed-size circular ring buffer $B \in \mathbb{R}^{W \times H \times 1}$, where $H = W = 160$ pixels. As new rows are scanned at a hardware scan speed $S = 10$ rows per clock tick, the buffer is updated by rolling the old rows up and inserting the new rows at the bottom:
\begin{equation}
B_{t} = \text{roll}(B_{t-1}, -S, \text{axis}=0)
\end{equation}
\begin{equation}
B_{t}[-S:, :, :] = \text{Ingested\_Rows}
\end{equation}
Inference is blocked until the buffer is at least 50\% filled (representing the scanning of 80 rows) to ensure minimum visual context.

\subsection{Model Architecture and Grayscale Adapter}
We deploy a MobileNetV2 backbone pretrained on ImageNet as our feature extractor. Because MobileNetV2 requires 3-channel RGB inputs, we prepended a grayscale-to-RGB adapter consisting of a $1 \times 1$ Convolution layer with 3 filters:
\begin{equation}
X_{\text{RGB}} = \text{Conv2D}_{1 \times 1}(X_{\text{Gray}}, \text{filters}=3)
\end{equation}
This adapter maps the single-channel buffer state directly to 3 channels without copying data in host memory. The base model features are pooled via Global Average Pooling and passed through a Dense layer with 256 units, a Dropout layer ($r = 0.4$), and a Softmax output layer mapping to 8 wafer defect classes: \textit{clean}, \textit{crack}, \textit{edge\_defect}, \textit{open}, \textit{other}, \textit{scratch}, \textit{short}, and \textit{spot}.

During training, we freeze the first 100 layers of the MobileNetV2 base and fine-tune the remaining layers using the Adam optimizer ($lr = 1\times 10^{-4}$) under a balanced class weight loss function to handle minor dataset imbalances.

\subsection{Confidence-Gated Early Exit and Entropy}
For every tick after the buffer is 50\% full, the system runs inference on the current buffer state. Let $P = [p_1, p_2, \dots, p_K]$ be the softmax output vector over $K=8$ classes. The predicted class is $c^* = \arg\max P$, and the confidence is $p_{c^*}$.

We define the normalized Shannon entropy $H$ of the prediction to measure uncertainty:
\begin{equation}
H = -\frac{1}{\ln K} \sum_{i=1}^{K} p_i \ln(p_i + \epsilon)
\end{equation}
where $\epsilon = 1 \times 10^{-9}$ is a stability constant. The decision logic at each tick is governed by the following rules:
\begin{itemize}
    \item If $p_{c^*} \ge 0.85$ and $H \le 0.35$:
    \begin{itemize}
        \item If $c^* \neq \text{clean} \rightarrow$ \textbf{REJECT} (Early exit)
        \item If $c^* = \text{clean} \rightarrow$ \textbf{ACCEPT} (Early exit)
    \end{itemize}
    \item If $p_{c^*} \ge 0.85$ and $H > 0.35 \rightarrow$ \textbf{DEFER} for human review
    \item If $p_{c^*} < 0.85 \rightarrow$ \textbf{CONTINUE} scanning next rows
\end{itemize}

\subsection{Monte Carlo (MC) Dropout Uncertainty Path}
For borderline cases where the fast path output triggers a deferral or requires formal safety validation, the system switches to the MC Dropout path. The inference graph is reconstructed by forcing the Keras dropout layer active (\texttt{training=True}). 

We tile a single input frame $X$ into a batch of $N=5$ identical copies and run a single forward pass:
\begin{equation}
\mathcal{P} = \text{Model}(X_{\text{batch}}, \text{training}=\text{True}) \in \mathbb{R}^{N \times K}
\end{equation}
We compute the mean prediction $\mu$ and prediction variance $\sigma^2$ across the $N$ stochastic runs:
\begin{equation}
\mu = \frac{1}{N} \sum_{j=1}^{N} \mathcal{P}_j, \quad \sigma^2 = \frac{1}{N} \sum_{j=1}^{N} (\mathcal{P}_j - \mu)^2
\end{equation}
If the mean confidence $\mu_{c^*} \ge 0.85$ and the variance $\sigma^2_{c^*} \le 0.05$, the system emits a deterministic decision; otherwise, the wafer is labeled as uncertain (\textbf{DEFER}).

\section{Experimental Evaluation}

\subsection{Dataset and Setup}
The model was trained on a dataset of 541 wafer surface defect patch images divided into 8 classes. The test dataset consists of 150 images (20 images for clean, crack, edge\_defect, other, scratch, spot; 15 images for open, short). Benchmarks were executed on an NVIDIA GeForce RTX 4050 Laptop GPU (6GB VRAM, Windows 11, CUDA 12.6, TensorRT 10.3, onnxruntime-gpu 1.19.2).

\subsection{Inference Latency Results}
We compared the inference latency of the exported ONNX model across three execution providers. TensorRT compiling with FP16 precision achieved significant speedup, as shown in Table~\ref{tab:inference_latency}.

\begin{table}[h]
\centering
\caption{Inference Latency Across Execution Providers}
\label{tab:inference_latency}
\begin{tabular}{lccccc}
\toprule
\textbf{Provider} & \textbf{Mean (ms)} & \textbf{Median (ms)} & \textbf{P95 (ms)} & \textbf{P99 (ms)} & \textbf{FPS} \\
\midrule
CPU Baseline & 2.64 & 2.46 & 3.12 & 3.72 & 379.1 \\
CUDA & 2.74 & 2.33 & 4.01 & 6.38 & 364.8 \\
\textbf{TensorRT FP16} & \textbf{0.79} & \textbf{0.80} & \textbf{1.05} & \textbf{1.17} & \textbf{1262.2} \\
\bottomrule
\end{tabular}
\end{table}

TensorRT FP16 achieves a mean inference time of \textbf{0.79~ms}, which is \textbf{3.34$\times$ faster} than the CPU baseline and represents a throughput of 1,262 frames per second.

\subsection{End-to-End Pipeline Latency}
We profiled the entire pipeline latency (from image availability to decision output) comparing a standard file-based pipeline with the OPTI-FAB stream pipeline under TensorRT FP16 acceleration.

\begin{table}[h]
\centering
\caption{End-to-End Pipeline Latency Comparison}
\label{tab:pipeline_latency}
\begin{tabular}{lrr}
\toprule
\textbf{Metric} & \textbf{File-Based Pipeline} & \textbf{OPTI-FAB Stream} \\
\midrule
Mean Latency & 8.30 ms & 3.00 ms \\
Median Latency & 7.74 ms & 2.47 ms \\
P95 Latency & 11.41 ms & 4.83 ms \\
Min Latency & 6.27 ms & 2.02 ms \\
Throughput & 120.5 fps & 333.3 fps \\
\bottomrule
\end{tabular}
\end{table}

OPTI-FAB reduces mean pipeline latency by \textbf{63.9\%} (from 8.30~ms to 3.00~ms), achieving a system speedup of **2.77$\times$**.

\subsection{Classification Performance}
OPTI-FAB achieved an overall accuracy of **90\%** on the 150 test samples. The per-class evaluation is detailed in Table~\ref{tab:classification_performance}.

\begin{table}[h]
\centering
\caption{Per-Class Evaluation Metrics}
\label{tab:classification_performance}
\begin{tabular}{lcccc}
\toprule
\textbf{Defect Class} & \textbf{Precision} & \textbf{Recall} & \textbf{F1-Score} & \textbf{Support} \\
\midrule
clean & 0.95 & 1.00 & 0.98 & 20 \\
crack & 0.95 & 1.00 & 0.98 & 20 \\
edge\_defect & 0.67 & 1.00 & 0.80 & 20 \\
open & 1.00 & 1.00 & 1.00 & 15 \\
other & 1.00 & 0.50 & 0.67 & 20 \\
scratch & 0.87 & 1.00 & 0.93 & 20 \\
short & 1.00 & 1.00 & 1.00 & 15 \\
spot & 1.00 & 0.75 & 0.86 & 20 \\
\midrule
\textbf{Accuracy} & & & \textbf{0.90} & \textbf{150} \\
\bottomrule
\end{tabular}
\end{table}

The model demonstrates excellent performance on critical structural defects, with perfect recalls on \textit{open}, \textit{short}, \textit{scratch}, and \textit{crack}. The lower recall on \textit{other} is primarily due to confusion with the \textit{edge\_defect} class, representing borderline aesthetic anomalies.

\section{Discussion and Factory Integration}
By reducing the pipeline latency to a mean of 3.00~ms, OPTI-FAB aligns perfectly with the target latency budget $T_{\text{budget}} = 3.0\text{ ms}$. At a transportation speed of 100~mm/s, the physical movement of the wafer during the inspection lifecycle is limited to $0.30\text{ mm}$ (compared to $0.83\text{ mm}$ for file-based pipelines), rendering inline sorting and sorting gates highly feasible on real fabrication floors.

Furthermore, because early exit triggers at 50\% to 68\% of frame scanning completion, the system saves up to 50\% of sensor scanning time on clear or clearly defective parts, allowing for immediate mechanical routing and increased production line throughput.

\section{Conclusion}
OPTI-FAB demonstrates that wafer quality control can be optimized as a real-time systems problem. By integrating stream-aware pixel processing, early exits via Shannon entropy, and edge GPU TensorRT compilation, the system breaks the file I/O bottleneck. Future research will explore hardware-in-the-loop validation on embedded boards (e.g., NXP i.MX8) and federated learning strategies to allow collaborative training across multiple fabs without exposing sensitive process data.

\bibliographystyle{IEEEtran}
\bibliography{references}

\end{document}
